\newpage
\section{Entwurf und Umsetzung des Ausführungsmodells} \label{sec:Entwurf}

Dieses Kapitel beantwortet U2. Es beginnt mit den Randbedingungen der Zielplattform, weil diese
den Entwurfsraum verengen, bevor eine einzelne Technologie geprüft wird. Darauf folgt die
Vorauswahl der Engine anhand der Muss-Kriterien und zuletzt der Aufbau des Modells in der
bestehenden Kubernetes-Umgebung.

\subsection{Randbedingungen der Zielplattform}
\label{sec:randbedingungen}

% GRUNDLAGEN: Kubernetes-Begriffe (Pod, Container, RuntimeClass) und cgroups v2
% gehoeren als Definition nach Kapitel 2. Hier bleibt nur ihre Folge fuer den Entwurf.

Drei Eigenschaften der bestehenden Umgebung stehen nicht zur Disposition.

Die erste ist das Fehlen eines Hypervisors. Firecracker setzt Lese- und Schreibzugriff auf \texttt{/dev/kvm}
voraus \autocite{the_firecracker_authors_getting_2026}, einen Zugriff, den der Zielcluster nicht
bereitstellt. MicroVM-basierte Ansätze scheiden damit aus, bevor ihre Isolationseigenschaften
bewertet werden. Der Ausschluss ist ein Plattformausschluss und kein Sicherheitsurteil, denn die
von \textcite{agache_firecracker_2020} beschriebene Grenze ist die stärkste der betrachteten.

Die zweite ist die bereits eingesetzte Sandbox-Laufzeit. M3 untersagt, die dem Gastcode zugängliche Oberfläche des Host-Kernels
gegenüber dem Ausgangszustand zu vergrößern. Die äußere Grenze ist damit gesetzt und keine
Entwurfsvariable mehr. Ein Modell, das diese Schicht zugunsten kürzerer Startzeiten aufgäbe,
verletzte M3, auch wenn es die Soll-Kriterien besser bediente. Der Aufwand der Schicht ist bekannt
und in die Entscheidung des Ausgangszustands bereits eingegangen \autocite{young_true_2019}. Für
das Bedrohungsmodell folgt daraus, dass ein Ausbruch aus der V8-Sandbox in einer gVisor-Sandbox
endet und nicht auf dem Host-Kernel. Hardwareseitige Seitenkanäle deckt gVisor nach eigener
Darstellung nicht ab \autocite{the_gvisor_authors_security_2026}, und der Sentry gehört nach
\autoref{sec:angreifermodell} selbst zur \ac{TCB}.

Die dritte betrifft die Granularität, auf der sich Ressourcen überhaupt begrenzen lassen. Kubernetes vergibt Ressourcenkontingente
je Container \autocite{the_kubernetes_authors_resource_2026} und kennt keine Schnittstelle, um
Prozessen innerhalb eines Containers getrennte Kontingente zuzuweisen. Auch der Thread-Modus von
cgroup v2 hilft nicht weiter, weil dort nur die Controller für Rechenzeit, Prozessorzuordnung,
Leistungsereignisse und Prozesszahl wirken, nicht aber der Speicher-Controller
\autocite{the_kernel_development_community_control_2026}. Der naheliegende Ausweg, je Aufgabe
einen eigenen Container im selben Pod zu führen, trägt unter der gVisor-Laufzeitklasse ebenfalls
nicht. Dort wird das Kontingent allein für den Sandbox-Prozess gesetzt, während die je Container
deklarierten Werte innerhalb der Sandbox nicht durchgesetzt werden
\autocite{googlegvisor_contributors_container_2019}.

Unabhängig davon wäre dieser Ausweg auch seinem Zweck nach der falsche. Container eines Pods
teilen sich Adresse und Portraum, finden einander über \texttt{localhost} und können über die
üblichen Mechanismen der Interprozesskommunikation miteinander sprechen. Die Dokumentation
empfiehlt das Muster ausdrücklich nur für eng gekoppelte Container
\autocite{the_kubernetes_authors_pods_2026}. Mehrere Container in einem Pod sind ein Mittel der
Zusammenarbeit und keines der Trennung. Der Pod ist damit zugleich die kleinste Einheit, auf der
ein vom Kern durchgesetztes Speicherkontingent existiert, und die kleinste, die Kubernetes als
Isolationseinheit anbietet.

\subsection{Vorauswahl der Engine}
\label{sec:vorauswahl}

Dieser Abschnitt legt fest, welche JavaScript-Engine das Modell ausführt. Entschieden wird allein
über die Muss-Kriterien M1 bis M5 und, wo mehrere Kandidaten dieselben Kriterien erfüllen, über
einen Dominanzvergleich, der ohne Gewichtung auskommt. Für Aussagen darüber, was eine
Laufzeitumgebung zusichert, wird deren normative Herstellerdokumentation herangezogen, denn für
einen Ausschluss ist die dokumentierte Zusicherung aussagekräftiger als eine Einzelmessung.

\subsubsection{Äußere und innere Isolationsgrenze}
\label{sec:zwei-grenzen}

% GRUNDLAGEN: Die Begriffe Isolate und Context sowie ihr Verhaeltnis zum Heap gehoeren
% als Definition nach Kapitel 2.

Ein verbreiteter Vergleich stellt Container und sprachbasierte Isolates als Alternativen
gegenüber. Diese Gegenüberstellung trägt nicht, weil beide an unterschiedlichen Stellen wirken.
Ein Container ist eine äußere Grenze, die vom Betriebssystemkern durchgesetzt wird, ein Isolate
eine innere Grenze innerhalb eines Prozesses. Eine äußere Grenze existiert in jedem Entwurf, auch
dann, wenn sie nicht eigens gewählt wird. Offen sind daher zwei andere Fragen, nämlich ob
zusätzlich eine innere Grenze eingezogen wird und wie viele Ausführungseinheiten sich eine Grenze
teilen.

Für M1 lässt sich das auflösen. Die Praxis der Anbieter zieht die Mandantengrenze durchgehend
außen: \textcite{reis_site_2019} führten die Prozesstrennung in Chrome ein, weil die
sprachbasierte Trennung mehrerer Origins nicht zuverlässig durchsetzbar blieb,
\textcite{cloudflare_security_2026} ergänzt die eigene Isolate-Plattform um eine Prozesssandbox
mit Namespaces und \texttt{seccomp}, und \textcite{deno_land_inc_how_2026} beschreibt für die
eigene mandantenfähige Plattform einen Prozess je Deployment, Namespaces, seccomp, cgroups und
einen Watchdog. Kein Kandidat erfüllt M1 also über die innere Grenze. Da die Fallstudie jedem
Mandanten bereits einen eigenen Execution Service gibt, erfüllt sie M1 im Ausgangszustand, und
das Kriterium scheidet keinen Kandidaten aus. Es legt fest, dass die mandantenweise Grenze
erhalten bleibt, und verschiebt die Entwurfsfrage auf die innere Grenze.

\subsubsection{Anwendung der Muss-Kriterien}
\label{sec:ausschluesse}

Vier Bauarten scheiden hier aus, drei an den Muss-Kriterien und eine am Gegenstand der Arbeit.

Node.js selbst kann als Vertrauensgrenze M4 nicht erfüllen. Das Bedrohungsmodell des Node.js-Projekts erklärt den ausgeführten Code ausdrücklich für
vertrauenswürdig, einschließlich dynamisch geladenen Codes, und schließt daraus, dass ein
Szenario, das ein bösartiges Modul voraussetzt, keine Schwachstelle in Node.js darstellen kann
\autocite{the_nodejs_project_security_2026}. Ein Skript, das \texttt{require} aufrufen kann,
erreicht Netz und Dateisystem an jedem vermittelten Binding vorbei, und ein Ausbruch aus einem
Worker wäre nach dieser Festlegung kein zu behebender Fehler. Das seit Version 20 verfügbare
Berechtigungsmodell bestätigt die Einordnung und bezeichnet sich selbst als \enquote{seat belt}
ohne Garantie gegenüber bösartigem Code \autocite{the_nodejs_project_permissions_2026}. Node.js
ist damit nicht schwach bewertet, sondern als Vertrauensgrenze nicht bewertbar. Als
Laufzeitumgebung innerhalb einer fremden Grenze bleibt es tragfähig, und in dieser Rolle erfüllt
der Ausgangszustand M1 und M3.

Ein Berechtigungsmodell an der Stelle einer Mandantengrenze scheitert an M1 und an M4. Die
Berechtigungsmodelle von Node.js und Deno wirken prozessweit und trennen keine Mandanten
innerhalb eines Prozesses, womit M1 nicht erfüllt ist. Für M4 nennt Deno selbst zwei Umgehungen.
Ein Skript, das einen neuen Deno-Prozess starten darf, kann diesen mit \texttt{-{}-allow-all}
starten, und eine über \ac{FFI} geladene native Bibliothek setzt Systemaufrufe direkt ab.
Dieselbe Dokumentation empfiehlt für nicht vertrauenswürdigen Code zusätzlich Sandboxing auf
Betriebssystemebene sowie gVisor oder Firecracker \autocite{the_deno_authors_security_2026}.

Bun scheitert an M2 und M4. Es liefert kein Berechtigungsmodell, der entsprechende Vorschlag ist ein offener Issue
\autocite{oven-shbun_contributors_bun_2026}. Unabhängig davon lässt es sich nicht als Bibliothek
einbetten, sodass ein Host keine Isolates erzeugen, begrenzen oder beenden könnte
\autocite{oven-shbun_contributors_feature_2026}. Ohne diese Kontrolle sind weder M2 noch M4
durchsetzbar.

WebAssembly als Mandantensprache scheitert an keinem Muss-Kriterium, sondern an der Oberfrage.
Diese fragt nach einem Ausführungsmodell für JavaScript-Funktionen. Ein Entwurf, in dem Mandanten
ihre Logik selbst nach WebAssembly übersetzen, beantwortet sie nicht und bricht zugleich den
gesamten Bestand aus \autoref{sec:bestandsanalyse}.

Nach diesen Ausschlüssen bleibt eine einzige Bauart übrig: eine in den Execution Service
eingebettete JavaScript-Engine, die eine eigene Isolate-Grenze je Ausführung zieht und Netz sowie
Dateisystem über vermittelte Bindings freigibt. Zwei technische Wege führen dorthin.

\subsubsection{V8-Isolate gegen eine nach WebAssembly übersetzte Engine}
\label{sec:wasm-zurueckgestellt}

% GRUNDLAGEN: WebAssembly als Ausfuehrungsformat, linearer Speicher und fehlende
% Codeerzeugung zur Laufzeit gehoeren als Definition nach Kapitel 2.

Der erste Weg bettet V8 unmittelbar ein und zieht die Isolate-Grenze nativ, mit
\ac{JIT}-Übersetzung. Der zweite übersetzt eine JavaScript-Engine selbst nach WebAssembly und
führt sie in einer Wasm-Sandbox aus. Die Mandanten schreiben in beiden Fällen JavaScript,
betroffen ist nur die Bauart der Engine. Beide erfüllen M1 bis M4 und lassen sich über die
Muss-Kriterien nicht trennen, wohl aber über einen Dominanzvergleich.

Beim Sicherheitsgewinn, dem einzigen Kriterium, an dem der zweite Weg ansetzen könnte, zeigt sich
kein Vorsprung. WebAssembly erlaubt keine Erzeugung und Ausführung von Maschinencode zur Laufzeit
\autocite{haas_bringing_2017}, was Angriffe auf erzeugten Code erschwert. Innerhalb des linearen
Speichers fehlen jedoch Schutzmechanismen, die native Programme besitzen, darunter \ac{ASLR},
Guard Pages und Stack Canaries \autocite{lehmann_everything_2020}, und Ausbrüche aus der
Wasm-Laufzeitumgebung selbst sind dokumentiert \autocite{bytecode_alliance_wasmtimes_2026}. Beide
Wege laufen ohnehin innerhalb derselben gVisor-Sandbox, sodass die zweite Sandbox-Schicht keine
Grenze ersetzt, sondern sie verdoppelt.

Auf den übrigen Kriterien unterliegt der zweite Weg. Ohne \ac{JIT} bleibt eine so übersetzte
Maschine interpretierend \autocite{igalia_five_2026}, was gegen S2 und S4 wirkt. Der Betrieb auf
Kubernetes erfordert zudem eigene containerd-Shims und damit einen zweiten Deployment-Stack neben
dem bestehenden, was gegen S5 wirkt. \textcite{shillaker_faasm_2020} belegen zwar einen
geringeren Ressourcenbedarf für serverloses Rechnen, und \textcite{mainas_sandboxing_2024}
übertragen den Befund auf mandantenfähige Umgebungen. Beide vergleichen jedoch gegen
Container-Isolation und nicht gegen ein eingebettetes, JIT-fähiges Isolate, sodass sich der
Befund auf den hier zu entscheidenden Vergleich nicht überträgt.

Damit erreicht der Wasm-Weg auf keinem Kriterium ein besseres Ergebnis als der V8-Weg und auf
mehreren ein schlechteres. Er scheidet bereits in der Vorauswahl aus. Eine quantitative
Nutzwertanalyse zwischen zwei derart ungleichen Kandidaten würde keinen zusätzlichen
Erkenntnisgewinn liefern. Sie ist deshalb nicht Gegenstand von \autoref{sec:nutzwertanalyse}.

\subsubsection{Die Kompatibilitätsschranke}
\label{sec:kompatibilitaet}

Die Kompatibilität des Bestands ist kein Muss-Kriterium, sie geht als S6 und S7 gewichtet in die
Nutzwertanalyse ein und scheidet deshalb keinen Kandidaten aus. Ihre Struktur bestimmt aber, mit
welchem Aufwand der gewählte V8-Weg umzusetzen ist, und gehört deshalb bereits hierher.

\autoref{sec:bestandsanalyse} weist 713 Mandanten (41,5\,\%) auf dem Hard Floor aus. Nach
\autoref{tab:blockierende-familien} scheitern davon 648, also 90,9\,\%, an einem
Netzwerk-Binding und an nichts anderem. Für die Umsetzung ist damit die entscheidende Aussage
getroffen: Der Aufwand der inneren Grenze konzentriert sich auf eine einzige Fähigkeit, und er
verteilt sich im Verhältnis von etwa zehn zu eins auf eine \ac{HTTP}-Gruppe und eine deutlich
kleinere Socket-Gruppe.

Für den Entwurf folgt daraus, dass beide Gruppen über einen Shim nachgebildet werden, der die
vorhandene Netzwerk-\ac{API} des Bestands abbildet. Ziel ist, dass die bestehenden Skripte unverändert weiterlaufen.
Der Shim gibt daher, wie das Original, einen rohen Socket frei und prüft dabei Host und Port
gegen die je Mandant hinterlegte Erlaubnis, bevor die Verbindung nach außen geht. Die Vermittlung
nach M4 liegt damit auf Ebene des Verbindungsaufbaus, nicht auf Ebene des Anwendungsprotokolls.
\autoref{sec:orchestrierung} legt die Umsetzung im Einzelnen fest.

\subsubsection{Ergebnis der Vorauswahl}
\label{sec:vorauswahl-ergebnis}

\begin{table}[H]
\centering
\caption{Anwendung der Muss-Kriterien und des Dominanzvergleichs auf die Kandidaten}
\label{tab:kandidaten}
\small
\begin{tabular}{@{}p{0.30\textwidth}ccccc p{0.15\textwidth}@{}}
\hline
Kandidat & M1 & M2 & M3 & M4 & M5 & Ergebnis \\
\hline
Ausgangszustand, Node.js mit wiederverwendetem Prozess & ja & nein & ja & nein & ja &
  scheidet aus \\
Node.js, frischer Prozess je Ausführung & ja & ja & ja & nein & offen &
  scheidet aus \\
JavaScript in WebAssembly je Ausführung & ja & ja & ja & ja & offen &
  dominiert, scheidet aus \\
V8-Isolate je Ausführung mit vermittelten Bindings & ja & ja & ja & ja & offen &
  \textbf{gewählt} \\
\hline
\end{tabular}
\end{table}

Der Ausgangszustand verletzt M2, weil aufeinanderfolgende Ausführungen sich den Prozessspeicher
teilen, und M4, weil dem Skript die vollständige Node-Schnittstelle offensteht.

Die zweite Zeile ist für die Arbeit die aufschlussreichere. Ein frischer Node-Prozess je
Ausführung stellt Zustandsfreiheit her und erfüllt M2, scheitert aber weiterhin an M4.
Zustandsfreiheit und vermittelter Zugriff sind also zwei verschiedene Eigenschaften, und die
naheliegende, billige Änderung liefert nur die erste. Ein Skript mit vollständiger
Node-\ac{API} erreicht jeden erreichbaren Endpunkt, und ein Kontrollpunkt lässt sich nicht
einziehen, weil die Laufzeitumgebung den Code als vertrauenswürdig behandelt
\autocite{the_nodejs_project_security_2026}. Für zur Laufzeit erzeugte Skripte ist das nach V2 die
maßgebliche Lücke, denn ein über seine Eingaben beeinflusstes Skript kann Daten abfließen lassen,
auf die es ohnehin zugreifen darf \autocite{greshake_not_2023}.

% GRUNDLAGEN: Das Verhaeltnis von V8, rusty_v8 und deno_core gehoert als
% Einordnung nach Kapitel 2.
Offen bleibt, auf welcher Ebene die Engine eingebettet wird, denn das bestimmt den
Umsetzungsaufwand. Der Binding Layer \texttt{rusty\_v8} stellt Isolate, Context, Skript und
Scopes bereit, aber weder Event Loop noch Modulauflösung noch einen Mechanismus, über den der Host
Funktionen in das Isolate reicht. Darauf setzt \texttt{deno\_core} auf und liefert genau diese
Schicht mit, einschließlich der Ops, über die Host-Bindings registriert werden. Es bildet zugleich
den Kern der Deno-Laufzeitumgebung \autocite{deno_land_inc_how_2026}. Für M4 ist das mehr als eine
Bequemlichkeit, weil eine fertige Schicht für die Vermittlung die selbst zu schreibende und selbst
abzusichernde Fläche verkleinert. Die Einbettung erfolgt deshalb über \texttt{deno\_core}. Das
zugehörige Wartungsrisiko behandelt \autoref{sec:betriebsplanung}.

\subsection{Aufbau des Ausführungsmodells}
\label{sec:orchestrierung}

\subsubsection{Ebenen der Trennung}
\label{sec:ebenen}

Das Modell trennt auf vier Ebenen, die von außen nach innen kurzlebiger und billiger werden. Nur
zwei davon tragen eine Grenze.

\begin{table}[H]
\centering
\caption{Ebenen des Ausführungsmodells und ihre jeweilige Aufgabe}
\label{tab:ebenen}
\small
\renewcommand{\arraystretch}{1.1}
\begin{tabularx}{\textwidth}{@{}>{\RaggedRight\arraybackslash}p{0.13\textwidth}
  >{\RaggedRight\arraybackslash\hsize=1.08\hsize}X
  >{\RaggedRight\arraybackslash\hsize=0.92\hsize}X
  >{\centering\arraybackslash}p{0.12\textwidth}@{}}
  \hline
  Ebene & Zuschnitt & Aufgabe & Kriterium \\
  \hline
  Pod & vorgehalten, bei Zuweisung einem Mandanten zugeordnet, gVisor-Runtime-Klasse &
    trennt Mandanten, setzt die harte Ressourcengrenze, hält die Startzeit vom Anfragepfad
    fern & M1, M3 \\[0.6ex]
  Prozess & genau einer je Pod &
    fällt mit der Pod-Grenze zusammen, trägt keine eigene & \\[0.6ex]
  Thread & je Prozess, zusätzlich ein Watchdog &
    bricht laufende Ausführungen ab, trägt keine Trennung & \\[0.6ex]
  Isolate & je Ausführung, aus dem Laufzeit-Snapshot &
    stellt Zustandsfreiheit her, trägt die Fähigkeitsgrenze & M2, M4 \\
  \hline
\end{tabularx}
\end{table}

\begin{figure}[H]
  \centering
  \newcommand{\prozess}[4]{%
    \node[hdr] (#1h) at (#2,#3) {Prozess};
    \node[lbl] (#1l) at ([yshift=-9pt]#1h.south west)
      {\begin{tabular}{@{}l@{}}
         genau einer je Pod, Speichergrenze aus dem Pod-cgroup\\
         Kontext in OpState: Mandant #4
       \end{tabular}};
    \node[iso] (#1a) at ([shift={(0.575,-0.42)}]#1l.south west) {Isolate};
    \node[note, anchor=west] (#1an) at ([xshift=5pt]#1a.east)
      {Skript des Webhook-Profils, Bindings: fetch, Socket};
    \node[iso] (#1b) at ([yshift=-0.60cm]#1a) {Isolate};
    \node[note, anchor=west] (#1bn) at ([xshift=5pt]#1b.east)
      {Skript des Agentenprofils, Bindings: fetch};
    \begin{pgfonlayer}{lproc}
      \node[proc, fit=(#1h)(#1l)(#1a)(#1an)(#1b)(#1bn)] (#1) {};
    \end{pgfonlayer}}
  \begin{tikzpicture}[
    iso/.style  ={draw, rounded corners=1pt, fill=black!10, minimum width=1.15cm,
                  minimum height=0.44cm, inner sep=1pt, font=\scriptsize},
    proc/.style ={draw, rounded corners=2pt, fill=white, inner sep=6pt},
    hdr/.style  ={font=\scriptsize\bfseries, inner sep=0pt, anchor=north west},
    lbl/.style  ={font=\scriptsize, inner sep=0pt, anchor=north west},
    note/.style ={font=\scriptsize\itshape, inner sep=0pt},
  ]
    \prozess{pa}{0}{0}{A}
    \prozess{pb}{0}{-4.6}{B}

    \node[lbl] (podal) at ([shift={(-0.34,0.72)}]pa.north west)
      {Pod, dem Mandanten A zugeordnet, gVisor-Runtime-Klasse};
    \begin{pgfonlayer}{lpod}
      \node[draw, rounded corners=3pt, fill=black!4, inner sep=11pt,
            fit=(podal)(pa)] (poda) {};
    \end{pgfonlayer}

    \node[lbl] (podbl) at ([shift={(-0.34,0.72)}]pb.north west)
      {Pod, dem Mandanten B zugeordnet, gVisor-Runtime-Klasse};
    \begin{pgfonlayer}{lpod}
      \node[draw, rounded corners=3pt, fill=black!4, inner sep=11pt,
            fit=(podbl)(pb)] (podb) {};
    \end{pgfonlayer}

    \node[lbl] (nodel) at ([shift={(0.14,0.68)}]poda.north west)
      {Kubernetes-Knoten, Pods aus einem vorgehaltenen Pool};
    \begin{pgfonlayer}{lnode}
      \node[draw, dashed, rounded corners=3pt, inner sep=11pt,
            fit=(nodel)(poda)(podb)] {};
    \end{pgfonlayer}
  \end{tikzpicture}
  \caption{Verschachtelung der Ausführungsebenen am Beispiel zweier Mandanten}
  \label{fig:ebenen}
\end{figure}

% GRUNDLAGEN: Kaltstartverdeckung durch vorgehaltene Instanzen in serverlosen
% Plattformen gehoert als Verfahren nach Kapitel 2.

Der vorgehaltene Pool folgt dem von \textcite{lin_mitigating_2019} beschriebenen Vorgehen,
Kaltstarts durch vorgehaltene Instanzen zu verdecken, statt den Start selbst zu beschleunigen.
Entscheidend ist die Regel, dass das Nachziehen des Pools einer Vorratspolitik folgt und niemals
auf dem Anfragepfad liegt. Wird sie verletzt, trägt eine Anfrage wieder die volle Startzeit eines
Pods. Das Verhalten des Ausgangszustands, bei fehlender Last auf null zu skalieren
\autocite{the_knative_authors_configuring_2026}, wird damit aufgegeben und durch eine Untergrenze
ersetzt.

Ein Pod trägt die Ausführungen genau eines Mandanten, wie schon im Ausgangszustand, und zwei
Mandanten teilen sich weder eine gVisor-Sandbox noch ein cgroup noch einen Adressraum. Container
und Prozess fallen mit der Pod-Grenze zusammen, weil nach \autoref{sec:randbedingungen} unterhalb
des Pods keine durchgesetzte Grenze existiert.

Aus demselben Grund unterscheidet das Modell die beiden Lastprofile nicht auf der Ebene der
Topologie. Sie unterscheiden sich allein in der Menge ihrer Bindings, und die wird nach
\autoref{sec:faehigkeitsmodell} je Skriptversion abgeleitet und je Isolate registriert, nicht je
Prozess. Ein Skript des Agentenprofils erhält also auch dann kein Socket-Binding, wenn im selben
Prozess ein Isolate läuft, das eines besitzt. Was ein zusätzlicher Schnitt nach Profil gebracht
hätte, liegt unterhalb dieser Grenze: Ein Abbruch, der den Prozess beendet, trifft alle
gleichzeitig laufenden Ausführungen, und nach einem Ausbruch aus der Engine liegen die
Ausführungen beider Profile im selben Adressraum. Beides betrifft ausschließlich Ausführungen
desselben Mandanten und ist gegenüber dem Ausgangszustand kein Rückschritt.
\autoref{sec:betriebsplanung} führt den Schnitt als offene Maßnahme.

Die Bindung eines Isolates an einen Thread ist keine Entwurfsentscheidung, sondern eine
Eigenschaft von \texttt{deno\_core}. Eine Laufzeitinstanz ist nicht zwischen Threads versetzbar
und bleibt an ihren Erzeugerthread gebunden. Ein Thread kann mehrere Instanzen nebenläufig
betreiben, solange diese ein- und ausgabegebunden sind. Rechenintensive Ausführungen blockieren
dabei ihre Nachbarn auf demselben Thread, weshalb ein gesonderter Watchdog-Thread nötig ist, der
den Abbruch auslösen kann, ohne selbst auf die Event Loop angewiesen zu sein. Eine Trennung trägt
der Thread nicht, denn alle Threads eines Prozesses teilen Adressraum und cgroup und enden bei
dessen Abbruch gemeinsam.

Die Isolate-Grenze je Ausführung erfüllt M2. Eine schwächere Variante, bei der ein Isolate
bestehen bleibt und lediglich ein frischer Context erzeugt wird, wäre erheblich billiger, stellt
Zustandsfreiheit aber nicht her, weil mehrere Contexts innerhalb eines Isolates denselben Heap
verwenden.

\subsubsection{Fähigkeitsmodell und Mandantenkontext}
\label{sec:faehigkeitsmodell}

Ein Isolate erreicht die Außenwelt nur über diejenigen Bindings, die der Host bei seiner Erzeugung
registriert hat. Ohne registriertes Binding existiert kein Weg nach außen, auch kein fehlerhaft
abgesicherter. Die Menge der registrierten Bindings ist damit die Grenze selbst und nicht deren
Bewachung, womit M4 strukturell erfüllt ist.

Der Mandantenkontext, also Kennung, Profil, zulässige Ziele für ausgehende Verbindungen und
Kontingente, liegt in \texttt{OpState}, dem Zustandsobjekt, das der Host jedem Isolate zuordnet.
Die Bindings lesen ihn dort aus dem Host-Code. Im JavaScript-Heap existiert er nicht, weshalb ein
Skript ihn weder auslesen noch überschreiben noch vortäuschen kann. Gegenüber dem Ausgangszustand
ist das der wesentliche Unterschied, denn dort trägt das Skript selbst die Angaben, mit denen es
sich ausweist. Für V2 ist die Eigenschaft zentral: Ein über seine Eingaben beeinflusstes Skript
kann die Frage, in wessen Namen es handelt, nicht beantworten und damit auch nicht falsch
beantworten.

Welche Bindings eine Skriptversion erhält, wird nicht je Mandant konfiguriert, sondern aus dem
Bundling abgeleitet. Der Ausgangszustand fasst jedes Kundenskript samt Abhängigkeiten bereits mit
esbuild zu einer Datei zusammen. Dieser Schritt löst den vollständigen Abhängigkeitsgraphen auf
und weiß daher, welche Node-Schnittstellen eine Skriptversion erreicht. Daraus folgt die Menge der
zu registrierenden Bindings mechanisch. Eine Skriptversion, die keine Socket-Schnittstelle
einbindet, erhält kein Socket-Binding.

Diese Ableitung setzt allerdings voraus, dass der Abhängigkeitsgraph zu einem Zeitpunkt
feststeht, an dem der Code noch nicht beeinflusst sein kann. Für gespeicherte Skripte trifft das
zu, weil das Bundling beim Hinterlegen läuft. Für zur Laufzeit erzeugten Code trifft es nicht zu.
Dort schreibt derjenige die Importe, dessen Beeinflussbarkeit V2 beschreibt, sodass eine
Ableitung aus dem Code dem Angreifer die Rechtevergabe überlassen würde. Die Menge der Bindings
für das Agentenprofil wird deshalb nicht abgeleitet, sondern plattformseitig fest gesetzt und
enthält kein Socket-Binding. Was für gespeicherte Skripte die feinere Lösung ist, wäre hier die
falsche.

Die Granularität ist damit zweistufig. Die Sicherheitsgrenze bleibt grob und verläuft je Mandant,
weil sie einen Pod kostet. Die Rechtevergabe verläuft fein und je Skriptversion, soweit ein
Bundling stattfindet, und ist andernfalls je Profil fest. Zwei Eigenschaften dieses Vorgehens
gehören in die Bewertung. Bei gespeicherten Skripten tritt fehlende Kompatibilität als Fehler
beim Bundling auf und nennt die fehlende Schnittstelle, während sie sich bei erzeugtem Code erst
zur Laufzeit zeigt. Und die Auflösung entzieht sich dem Mandanten, weil dieser den
Bundling-Schritt nicht kontrolliert. Nicht erfasst wird dynamisch aus Zeichenketten
zusammengesetztes Nachladen, das statisch nicht auflösbar ist und nach
\autoref{sec:bestandsanalyse} drei Mandanten betrifft.

\subsubsection{Kompatibilitätsschicht}
\label{sec:kompatibilitaetsschicht}

Nach \autoref{sec:bestandsanalyse} scheitern 648 der 713 blockierten Mandanten an einem
Netzwerk-Binding, verteilt auf eine \ac{HTTP}-Gruppe und eine etwa zehnmal kleinere
Socket-Gruppe. Der Umsetzungsaufwand der inneren Grenze konzentriert sich damit auf eine einzige
Fähigkeit. Beide Gruppen werden über einen Shim nachgebildet, der die vorhandene
Netzwerkschnittstelle auf ein vermitteltes Binding abbildet, sodass die bestehenden Skripte
unverändert weiterlaufen. Ziel ist ausdrücklich nicht, eine Abstraktion etwa einer
\ac{SQL}-Verbindung anzubieten.

Die fertige Node-Kompatibilitätsschicht der Deno-Laufzeitumgebung einzubinden wäre die
naheliegende Abkürzung, wird aber nicht gewählt. Sie bringt neben den benötigten
Netzwerkschnittstellen auch Bausteine für Prozesse, Dateisystem und Betriebssystemzugriff mit und
stellt damit genau diejenige Umgebungsautorität wieder her, wegen der Node.js in
\autoref{sec:ausschluesse} als Vertrauensgrenze ausgeschieden ist. Abgesichert wäre sie allein
durch dasselbe prozessweite Berechtigungsmodell, dessen Umgehungen der Hersteller selbst
dokumentiert \autocite{the_deno_authors_security_2026}. Ein selbst geschriebener Shim deckt
dagegen nur die Oberfläche ab, die der Bestand tatsächlich benötigt.

Ob dieser Zuschnitt trägt, ist durch Architekturbetrachtung nicht zu klären. Die Machbarkeit
wurde deshalb vorab an der schwierigsten Stelle des Bestands geprüft, nämlich an
\texttt{tedious}, dem Treiber unterhalb von \texttt{mssql}. Dieser bindet nicht nur einen rohen
Socket ein, sondern stellt eine bestehende Verbindung mitten im Datenstrom auf \ac{TLS} um, indem
er den Handshake in die Paketrahmen des \ac{TDS}-Protokolls verpackt.

Beide Teile sind belegt. Der unveränderte Treiber verbindet sich über einen selbst geschriebenen
Socket-Shim mit einer echten Datenbankinstanz und führt eine Abfrage aus, wobei host-seitig ein
einziges neues Binding für den Verbindungsaufbau nötig war, weil \texttt{deno\_core} Lesen,
Schreiben und Schließen bereits generisch über beliebige Ressourcen bereitstellt. Die
\ac{TLS}-Umstellung gelingt über eine Zustandsmaschine ohne eigene Ein- und Ausgabe, die
ausschließlich über synchrone Bindings angesprochen wird und die Ver- und Entschlüsselung
übernimmt, während der Treiber die Paketrahmung weiterhin selbst verantwortet. Beide Ergebnisse
sind mehrfach reproduziert und gegen Negativkontrollen geprüft, \autoref{sec:Evaluation} führt sie
im Einzelnen aus.

Der Vermittlungspunkt liegt bei der Socket-Gruppe auf Ebene des Verbindungsaufbaus. Geprüft werden
Zielhost und Zielport gegen die je Mandant hinterlegte Erlaubnis, bevor die Verbindung zustande
kommt. Danach besteht keine Einsicht in den Datenstrom mehr, weil dieser ein mandanteneigenes
Protokoll trägt und nach der \ac{TLS}-Umstellung zusätzlich verschlüsselt ist. Für die
\ac{HTTP}-Gruppe liegt der Vermittlungspunkt dagegen auf Ebene der einzelnen Anfrage. Diese
Ungleichheit ist eine Eigenschaft des Entwurfs und keine Lücke in seiner Umsetzung, und sie gehört
als solche in die Bewertung von M4.

\subsubsection{Begrenzung der Ressourcen}
\label{sec:ressourcenbegrenzung}

Die Begrenzung erfolgt auf zwei Ebenen, weil keine der beiden für sich genommen trägt.

Auf Ebene des Isolates greifen drei Mechanismen. Eine Obergrenze für den JavaScript-Heap, die über
einen Callback durchgesetzt wird, sobald sie erreicht ist. Ein zählender Allokator für Speicher,
den V8 außerhalb dieses Heaps anfordert. Und ein Watchdog, der eine laufende Ausführung abbricht.

Jeder dieser Mechanismen hat eine benennbare Grenze, und diese Grenzen begründen die zweite Ebene.
Die Heap-Obergrenze wird an den Haltepunkten der Speicherbereinigung geprüft und ist damit eine
Näherung, keine erzwungene Schranke. Sie erfasst zudem nur den JavaScript-Heap, weshalb der
zählende Allokator für Binärpuffer außerhalb dieses Heaps nötig ist. Der Abbruch durch den
Watchdog wirkt an den Prüfpunkten der Engine und damit nicht innerhalb eines laufenden nativen
Aufrufs. Vollständigkeit entsteht bei allen dreien durch Aufzählung der abgedeckten Pfade, und ein
nicht aufgezählter Pfad bleibt offen.

Auf Ebene des Pods greift deshalb ein cgroup je Mandant. Es erfasst jeden Allokationspfad, auch
die nicht aufgezählten, und stellt damit die harte Schranke dar. Feiner ist sie nach
\autoref{sec:randbedingungen} nicht zu bekommen. Weil ein Pod genau einen Mandanten trägt, fällt
die feinste erreichbare Einheit mit der Mandantengrenze zusammen. Der Kern wählt sein Opfer bei
Überschreitung zwar nach eigener Bewertung aus, doch jede in Frage kommende Ausführung gehört
demselben Mandanten. Die Zuordnung zum verursachenden Mandanten bleibt damit erhalten, die
Zuordnung zur verursachenden Ausführung leistet die Ebene des Isolates. Dieselbe Aufteilung findet
sich bei \textcite{deno_land_inc_how_2026}, wo cgroups neben Namespaces und seccomp als eigene
Schicht geführt werden.

Die Aufgabenteilung ist damit festgelegt. Die Mechanismen auf Ebene des Isolates liefern die
Attribution und einen im Skript abfangbaren Fehler, also eine Reaktion, die den Prozess nicht
beendet. Das cgroup liefert die Garantie. Die Belastbarkeit dieser Aufteilung prüft
\autoref{sec:Evaluation} anhand der Prüffälle aus \autoref{tab:prueffaelle}.

\subsubsection{Zustellung und Zwischenspeicherung des Skripts}
\label{sec:zustellung}

% GRUNDLAGEN: Snapshots als Verfahren, den Zustand einer initialisierten Instanz
% festzuhalten, gehoeren als Begriff nach Kapitel 2.

Der Ausgangszustand reicht das auszuführende Skript nach \autoref{sec:systemkontext} je Aufruf
über eine vorsignierte URL herein. Das Modell übernimmt dieses Vorgehen unverändert. Die
Alternative, einen Pod an eine Skriptversion zu binden, würde die Zahl der vorzuhaltenden Pods an
die Zahl der Skriptversionen koppeln und damit die Kaltstartverdeckung aus \autoref{sec:ebenen}
unbezahlbar machen.

Daraus folgt eine Unterscheidung, die für die Startzeit entscheidend ist. Vorbereitet werden kann
nur, was über Ausführungen hinweg gleich bleibt, und das hängt an der Herkunft des Skripts.

\begin{table}[H]
  \caption{Was sich je Herkunft des Skripts vorbereiten lässt}
  \label{tab:vorbereitung}
  \small
  \begin{tabularx}{\textwidth}{@{}>{\RaggedRight\arraybackslash}p{0.24\textwidth}
    >{\RaggedRight\arraybackslash}X >{\RaggedRight\arraybackslash}X@{}}
    \hline
     & gespeichertes Skript & zur Laufzeit erzeugtes Skript \\
    \hline
    Zustellung & vorsignierte URL je Aufruf & mit dem Aufruf \\[2pt]
    Bundling vorab & ja, beim Hinterlegen & nein \\[2pt]
    Menge der Bindings & aus dem Bundle abgeleitet & fest je Profil \\[2pt]
    Laufzeit-Snapshot & ja & ja \\[2pt]
    Skriptbezogenes Zwischenkompilat & möglich, je Skriptversion & nicht möglich \\[2pt]
    Aufnahme in den Zwischenspeicher & ja, über den Inhalts-Hash & nein, keine stabile Kennung \\
    \hline
  \end{tabularx}
\end{table}

Der \emph{Laufzeit-Snapshot} ist der Teil, der für jede Ausführung trägt. Er hält den Zustand
einer Instanz fest, in der die Erweiterungen der Einbettung und der JavaScript-Anteil der Shims
aus \autoref{sec:kompatibilitaetsschicht} bereits ausgewertet sind, und er entsteht einmalig beim
Bau der Binärdatei. Weil er kein Mandantenskript enthält, ist er für alle Mandanten, beide
Profile und jedes Skript derselbe. Das Verfahren entspricht dem von
\textcite{ustiugov_benchmarking_2021} untersuchten Vorgehen, den Zustand einer initialisierten
Instanz festzuhalten und für spätere Starts wiederzuverwenden.

Ein \emph{skriptbezogenes Zwischenkompilat} spart zusätzlich das wiederholte Übersetzen desselben
Quelltexts. Es setzt voraus, dass derselbe Quelltext mehrfach ausgeführt wird, und ist deshalb
allein für gespeicherte Skripte sinnvoll. Für ein Skript, das genau einmal läuft, gibt es nichts
wiederzuverwenden. Zwischengespeichert wird dabei ausschließlich Unveränderliches, also Quelltext
und übersetzter Code, niemals Heap-Zustand einer vorangegangenen Ausführung. Diese Abgrenzung ist
nicht nebensächlich, denn ein Snapshot des Mandantenskripts würde Heap-Zustand über Ausführungen
hinweg tragen und damit M2 verletzen. \autoref{tab:prueffaelle} führt den Fall als T3, der
dadurch auch den Laufzeit-Snapshot betrifft: Verändert ein Skript einen eingebauten Prototyp, darf
diese Änderung die nächste Ausführung nicht erreichen.

Nach \autoref{sec:bestandsanalyse} misst ein Bundle im Median 69,3\,KB und im 95.\ Perzentil
429,0\,KB. Schwer wiegt dabei nicht die Übertragung dieser Menge, sondern der konstante Anteil
eines Abrufs aus dem Objektspeicher, und der fällt bei jedem Aufruf erneut an. Der Prozess hält
deshalb einen Zwischenspeicher, dessen Schlüssel der Inhalt selbst ist. Vier Festlegungen
bestimmen ihn.

\emph{Er liegt im Speicher und nicht auf der Platte.} Für den Median-Mandanten mit zwei Versionen
sind das wenige hundert Kilobyte, im 95.\ Perzentil wenige Megabyte. Ein Zwischenspeicher auf der
Platte erzeugte dagegen bei jedem Zugriff Dateisystemaufrufe, die unter gVisor der Sentry
beantwortet, also diejenige Komponente, die nach \autoref{sec:randbedingungen} selbst zur \ac{TCB}
gehört. Für Datenmengen dieser Ordnung ist das ein unnötiger Umweg über eine Vertrauenskomponente.

\emph{Die Inhaltskennung kommt mit dem Aufruf.} Könnte der Prozess sie erst nach dem Laden
bestimmen, wäre die Prüfung wertlos, weil das Laden ja gerade vermieden werden soll. Der Aufruf
trägt deshalb einen Hash des Bundles, und der Prozess prüft die geladenen Bytes gegen ihn, bevor
er sie aufnimmt.

\emph{Aufgenommen wird nur, was eine stabile Kennung hat.} Zur Laufzeit erzeugter Code hat keine.
Ein Eintrag dafür wäre nie ein Treffer, und ohne diese Regel könnte ein Mandant über
Agentenaufrufe beliebig viele einmalige Skripte einlagern und damit die Bundles seines eigenen
Webhook-Pfads verdrängen. Das ist dieselbe Wirkung, die \autoref{sec:angreifermodell} unter A1 als
Erschöpfung einer geteilten Ressource beschreibt, hier auf den Zwischenspeicher angewandt.

\emph{Die Grenze ist ein Budget in Byte.} Einzelne Mandanten halten bis zu 151 gebündelte
Versionen, eine Grenze über deren Anzahl trüge also nicht. Das Budget zählt gegen das
Speicherkontingent des Pods aus \autoref{sec:ressourcenbegrenzung} und geht von dem ab, was den
Ausführungen zur Verfügung steht.

Der Zwischenspeicher nimmt dem zweiten Aufruf den Abruf ab, dem ersten nicht. Er wird deshalb
bereits bei der Zuweisung eines Pods an einen Mandanten mit denjenigen Bundles gefüllt, die dieser
Mandant gegenwärtig aufruft. Das ist kein zweiter Mechanismus, sondern derselbe zu einem früheren
Zeitpunkt, und er folgt derselben Regel wie der Pod-Pool aus \autoref{sec:ebenen}: Was sich
vorbereiten lässt, gehört nicht auf den Anfragepfad. Bei zwei Versionen im Median kostet das
wenige hundert Kilobyte und verschwindet neben den Sekunden, die ein Pod-Start ohnehin braucht.
Für das Agentenprofil greift es nicht, weil dort jedes Skript neu ist.

Ein gemeinsamer Zwischenspeicher für alle Pods eines Knotens würde mehr sparen, wird aber
verworfen. Er wäre von den Pods mehrerer Mandanten erreichbar und verriete über seine Antwortzeit,
ob ein anderer Mandant dasselbe Bundle hält. Dass solche Treffer vorkämen, belegt die
Bestandsanalyse, in der dieselbe Skriptkennung bei mehreren Mandanten auftritt.

Dass ein frisches Isolate je Ausführung das Bundle jedes Mal erneut auswerten muss, ist gegenüber
dem Ausgangszustand keine Verschlechterung, denn dort geschieht bereits dasselbe. Das Skript wird nach \autoref{sec:systemkontext} je Ausführung neu geladen, und sein
Eintrag im Modul-Zwischenspeicher wird verworfen, sodass das Übersetzen ebenfalls jedes Mal
anfällt. Bei einem Bundle von im Median 69,3\,KB zahlen beide Modelle diesen Preis.

Der Unterschied liegt in der Gegenleistung. Der Ausgangszustand erreicht damit, dass Zustand auf
Modulebene nicht in die nächste Ausführung gelangt, und setzt zusätzlich das temporäre Verzeichnis
zurück. Globale Variablen, veränderte Prototypen, registrierte Callbacks und offene Verbindungen
überdauern dagegen weiterhin. Das frische Isolate verwirft für denselben Aufwand den gesamten Heap und
erfasst damit jeden im Speicher liegenden Träger, den M2 benennt.

Vorladen und skriptbezogenes Zwischenkompilat sind damit keine Kompensation für eine selbst
eingeführte Last, sondern Verbesserungen gegenüber dem Ausgangszustand an einer Stelle, an der
dieser bisher bei jeder Ausführung voll bezahlt. Wie groß sie ausfallen, misst
\autoref{sec:Evaluation} an echten Bundles des Bestands und nicht an einem künstlichen Handler.

\subsubsection{Lebenszyklus einer Ausführung}
\label{sec:lebenszyklus}

Eine Ausführung läuft damit so ab. Die Zuteilung sucht den Pod des Mandanten und weist andernfalls
einen Pod des Pools zu. Der Prozess schlägt die mit dem Aufruf übergebene Inhaltskennung im
Zwischenspeicher nach und lädt das Bundle nur bei einem Fehltreffer über die vorsignierte URL,
prüft es gegen die Kennung und nimmt es auf, soweit Herkunft und Budget das zulassen. Aus dem
Laufzeit-Snapshot entsteht ein frisches Isolate, der Host registriert die für dieses Skript
vorgesehenen Bindings und legt den Mandantenkontext in \texttt{OpState} ab, dann wird ausgewertet
und der Einstiegspunkt aufgerufen. Anschließend wird das Isolate verworfen, der Eintrag im
Zwischenspeicher bleibt.

Für das Ende einer Ausführung gilt eine Festlegung, die aus der Vorstudie folgt. Eine Ausführung
endet nicht, wenn das Skript zurückkehrt, sondern wenn die Event Loop des Isolates leer ist. Eine
noch ausstehende Leseoperation hält eine eigene Referenz auf ihre Ressource, sodass das Schließen
der Ressource weder den Socket verwirft noch die Leseoperation beendet. Jede Ressource führt
deshalb eine Abbruchmöglichkeit mit, die beim Schließen ausgelöst wird, und jede im Hintergrund
laufende Operation wird so gekennzeichnet, dass sie das Isolate nicht am Leben hält. Ohne diese
Festlegung trüge ein Skript mit hängender Leseoperation seine Ausführung über deren Ende hinaus,
was T3 und T5 unmittelbar berührt.

\subsection{Betrieb und offene Maßnahmen}
\label{sec:betriebsplanung}

Für M5 ist maßgeblich, dass die Zuteilung beim Entzug eines Knotens laufende Ausführungen
abschließen lässt und neue auf verbleibende Pods leitet. Weil eine Ausführung kurz ist und keinen
Zustand über ihr Ende hinaus trägt, ist ein Abbruch verkraftbar und die Wiederholung einer
abgebrochenen Ausführung unbedenklich. Die Untergrenze des Pools ist so zu wählen, dass der Entzug
eines Knotens nicht dazu führt, dass Anfragen auf das Nachziehen warten.

Das in \autoref{sec:vorauswahl-ergebnis} benannte Wartungsrisiko betrifft die Einbettung selbst.
Die Bausteine von \texttt{deno\_core} erscheinen als
Vorabversionen in kurzen Abständen, was für ein Produkt mit langen Wartungszyklen eine dauerhafte
Bindung bedeutet. Der Entwurf begrenzt sie auf zwei Wegen. Übernommen wird nur die Kernschicht und
nicht die darauf aufbauende Kompatibilitätsschicht, wodurch die Zahl der nachzuführenden Bausteine
erheblich kleiner ausfällt. Und die selbst geschriebenen Shims bleiben unter eigener Kontrolle und
sind von Änderungen nur dort betroffen, wo diese ihre Schnittstelle ändern.

Eine zweite offene Maßnahme betrifft die Härtung der Engine. Die von den V8-Entwicklern
bereitgestellte Speichersandbox ist
über den regulären Auslieferungsweg von \texttt{rusty\_v8} derzeit nicht erreichbar, sie erfordert
eine vollständige Übersetzung der Engine aus dem Quelltext. Für den Prototyp wird deshalb nur
Pointer Compression aktiviert, die als vorgefertigtes Artefakt vorliegt und einen Teil der
Eindämmung leistet, indem sie Zeiger innerhalb des Heaps auf einen begrenzten Adressbereich
festlegt. Die vollständige Sandbox bleibt offen. Sie ist für M1 bis M5 nicht erforderlich, weil
die Mandantentrennung über die Pod-Grenze und nicht über die Engine läuft.

Der in \autoref{sec:ebenen} verworfene Schnitt nach Lastprofil schließlich wartet nicht auf eine
Änderung der Plattform. Ein zweiter Container im selben Pod wäre nach
\autoref{sec:randbedingungen} auch dann das falsche Mittel, wenn die Laufzeitklasse Kontingente je
Container durchsetzte. Wer den Schnitt will, braucht einen zweiten Pod je Mandant und damit die
doppelte Vorhaltung. Es ist eine Abwägung zwischen einer nur mandantenintern wirkenden Eindämmung
und S1 sowie S3, und keine technische Schranke.
